% TOP_PROBLEM: {"id":30007710,"problem_number":"LOCAL-30007710","title":"Adaptive experimental inference","category_id":21,"category":{"id":21,"name":"economics","display_name":"Economics","description":"Open economic theory statements, published research agendas, and proposed scoped specifications; question provenance and historical priority are qualified per record.","slug":"economics","order_index":21,"created_at":"2026-10-08T00:00:00Z"},"status":"provenance_pending","external_url":null,"legacy_ids":[],"published":false,"statement":"Characterize regret and simultaneous-inference trade-offs for adaptive experiments with bounded outcome drift and minimum exploration.","economics":{"original_id":199,"field":"Econometrics, causal inference and measurement","kind":"design","formalization_basis":"proposed_specification","source_status":"not_verified","priority_status":"not_established","priority_note":"No exact open-problem passage was verified. The supporting publication does not establish priority or unresolved status.","earliest_verified_reference":null,"current_reference":{"authors":["Aurélien Bibaut","Nathan Kallus"],"title":"Demystifying Inference After Adaptive Experiments","year":2024,"url":"https://arxiv.org/pdf/2405.01281","doi":"10.48550/arXiv.2405.01281","locator":"Sections 2–5; 2024 preprint, published 2025","evidence_summary":"Reviews solved adaptive-inference methods; no explicit statement of the proposed nonstationary joint welfare–inference frontier was verified."},"formalization_note":"The cited primary work supports the subject or supplies solutions in particular settings, but no exact open-question passage for this specification was verified. The mathematical target and restrictions are proposed here, with no canonical-conjecture or priority claim. Corrects the original catalogue's author attribution to Aurélien Bibaut and Nathan Kallus. Independent math review tightened feasibility, existence, or estimand assumptions where required."},"statusQualification":"Provenance pending: an explicit published statement of this exact open question was not verified. This is a catalogue-authored candidate specification, not a certified open conjecture. The cited primary work supports the subject or supplies solutions in particular settings, but no exact open-question passage for this specification was verified. The mathematical target and restrictions are proposed here, with no canonical-conjecture or priority claim. Corrects the original catalogue's author attribution to Aurélien Bibaut and Nathan Kallus. Independent math review tightened feasibility, existence, or estimand assumptions where required.","statusReviewedAt":"2026-10-08"}
% FIRST_POSTED: 2026-10-08
% ENTRY_KIND: statement_only
% !TeX program = lualatex
\documentclass[11pt]{article}
\usepackage{fontspec}
\usepackage[a4paper,margin=25mm]{geometry}
\usepackage{amsmath,amssymb,mathtools}
\usepackage{xurl}
\usepackage[hidelinks]{hyperref}
\newcommand{\sref}[2]{\hyperlink{source:#1:#2}{[S#2]}}
\setlength{\emergencystretch}{3em}
\begin{document}
% problemId:30007710
\section*{Adaptive experimental inference}\label{30007710}
\subsection{Definitions and mathematical statement}
Consider a sequential experiment over \(T\) participants with two treatments \(a\in\{0,1\}\) and bounded potential outcomes \(Y_t(a)\in[0,1]\). Conditional mean rewards \(\mu_{a,t}=E[Y_t(a)\mid\mathcal H_{t-1}]\) may change with time, where \(\mathcal H_{t-1}\) is recorded history. Assume conditional randomization given history and an almost-sure drift budget \(\sum_{t=2}^{T}\max_a|\mu_{a,t}-\mu_{a,t-1}|\leq V\). Assignment probability \(p_t\), measurable from history, satisfies \(\epsilon\leq p_t\leq1-\epsilon\) for fixed \(0<\epsilon<1/2\). Define the running average effect \(\tau_t=t^{-1}\sum_{s=1}^{t}(\mu_{1,s}-\mu_{0,s})\) and welfare regret \(R_T=\sum_{t=1}^{T}\{\max_a\mu_{a,t}-\mu_{A_t,t}\}\), where \(A_t\) is the assigned treatment. Develop assignment and inference rules whose confidence intervals \(C_t\) satisfy \(\Pr(\tau_t\in C_t\text{ for every }t\leq T)\geq1-\alpha\), for stated \(\alpha\in(0,1)\), uniformly over the bounded-outcome drift class. Characterize attainable trade-offs between expected regret and interval width as functions of \(T,V,\epsilon\), and compare them with fixed randomization. Regret compares one-step choices along realized history, rather than alternative state trajectories. Fixed positive exploration can imply linear regret even without drift. Arbitrary unbounded drift is excluded explicitly. Adaptive inference already has established solutions; the nonstationary joint welfare-and-inference optimization is a proposed extension, and the supporting reference does not verify its exact open status. Outcomes are assessed in a prespecified participant population with consistent measurement across the observation horizon.

\par\textbf{Formulation and scope.} The cited primary work supports the subject or supplies solutions in particular settings, but no exact open-question passage for this specification was verified. The mathematical target and restrictions are proposed here, with no canonical-conjecture or priority claim. Corrects the original catalogue's author attribution to Aurélien Bibaut and Nathan Kallus. Independent math review tightened feasibility, existence, or estimand assumptions where required.

\par\textbf{Open-question provenance.} An explicit published open-question statement was not verified. Sources below are supporting literature and must not be treated as a first listing.

\par\textbf{Historical priority.} No exact open-problem passage was verified. The supporting publication does not establish priority or unresolved status.

\subsection{Short English statement}
Characterize regret and simultaneous-inference trade-offs for adaptive experiments with bounded outcome drift and minimum exploration.

\subsection{Sources}
\begin{itemize}\raggedright
\item[\textnormal{[S1]}]\hypertarget{source:30007710:1}{} Aurélien Bibaut, Nathan Kallus. 2024. Demystifying Inference After Adaptive Experiments. Sections 2–5; 2024 preprint, published 2025.
\url{https://arxiv.org/pdf/2405.01281}.
DOI: 10.48550/arXiv.2405.01281.
{\small\textit{Supporting or current literature; see evidence qualification. Reviews solved adaptive-inference methods; no explicit statement of the proposed nonstationary joint welfare–inference frontier was verified.}}
\end{itemize}
\end{document}
